\section{Recurrent two-ended packets and the full radial core}
\label{sec:v96-packets}
The occupation theorem estimates a positive packet itself, rather than
only a difference between its two traces. The estimated family can
contain arbitrarily short exact-event segments: the anchors need not
lie in the original segment, but must return to the same selected event
on the same regular physical word.

Set $\alpha=129/128$ and $\gamma=1/32$, and write
\[
 H_{a,\delta}=\Pi\cap\{z<a\ \hbox{or}\ \ell<\delta\}.
\]
The intersection is counted once. Its unweighted density satisfies
$U_{H_{a,\delta}}\le p_a+e^-_\delta$, so the retained critical-volume
and backward-end estimates give, at fixed auxiliary band,
\begin{align}
 m^2\int_JU_{H_{a,\delta}}
 \le{}&C(a^\alpha+\delta^\gamma)(|J|+B_*^{-1})\notag\\
 &+C_{B_*}(1+B_*|J|)
    \{m^2a^{1/2}\theta_{B_*}^m
       +m^3\min(\theta_{B_*}^m,\sqrt\delta)\}.
                                                \label{eq:v96-union-finite}
\end{align}
This is an estimate on the original exact event, not on the enlarged
comparison domain.

\begin{theorem}[Untruncated recurrent-source height]
\label{thm:v96-source-height}
For fixed $r,D,\beta,a,\delta>0$ in their stated ranges,
\begin{equation}\label{eq:v96-union-height}
 \limsup_m\sup_{R,n,k}m^2
 \|R_{H_{a,\delta}^{\rm adm}}^\varepsilon\|_\infty
       \le C\frac{D}{\beta}(a^\alpha+\delta^\gamma).
\end{equation}
The constant is independent of $\varepsilon$ and of the previously
fixed comparison radius. At finite count the right side before taking
limits is $D/(\beta r)$ times the right side of
\eqref{eq:v96-union-finite} with $|J|=2r$.

The same conclusions separately hold for $H=\Pi\cap\{z<a\}$
with only $a^\alpha$, and for $H=\Pi\cap\{\ell<\delta\}$
with only $\delta^\gamma$. In particular, for fixed $D,\beta$,
the original choice $a=h^{1/8}$, $\delta=h^{1/4}$ gives
\begin{equation}\label{eq:v96-cofinal-power}
 \mathcal H_M(R_{H_{a,\delta}^{\rm adm}}^\varepsilon)
                                    \le C(D/\beta)h^{1/128}.
\end{equation}
The radius may be a prescribed function of $h$, fixed before each
collision limit.
\end{theorem}
\begin{proof}
Apply \eqref{eq:v96-unweighted-anchor} and then
\eqref{eq:v96-union-finite} on $[t-r,t+r]$. The estimate is uniform
in the interval location. Fix all displayed parameters and $B_*$,
take the collision limsup, and only then enlarge $B_*$. The floor
$B_*^{-1}/r$ disappears in this order, and both spectral remainders
vanish. The leading interval length cancels the $r^{-1}$ factor.
The two single restrictions use the corresponding single input.
Finally $a^\alpha=h^{129/1024}$ and
$\delta^\gamma=h^{1/128}$, proving \eqref{eq:v96-cofinal-power}.
\end{proof}

This radial conclusion contains every annulus below $a$ at finite
count, including annuli whose locations depend on $m$. It is not a
sum of fixed-annulus limsups. It does not assert that every such
annulus is admitted: comparison survival, physical compression and
same-label occupation remain explicit tests.

\begin{corollary}[No additional radial-return test in the inner core]
\label{cor:v96-inner-core}
Let $0<r\le a^2/(8\Lambda_\circ)$. Targets with $z(t)<a/2$
which pass the continuation tests for $H=\Pi$ also pass them for
$H=\Pi\cap\{z<a\}$. Their complete weighted source has ordered
height at most $C(D/\beta)a^{129/128}$, with no lower gradient cutoff.
\end{corollary}
\begin{proof}
On any available backward interval $z$ decreases. On an available
forward interval, $z(t+s)^2\le a^2/4+2\Lambda_\circ r\le a^2/2$.
Thus every sample returning to $\Pi$ also lies in $z<a$. Both
occupation fractions coincide, and the distortion test is unchanged.
Apply the single-critical-restriction part of
Theorem~\ref{thm:v96-source-height}.
\end{proof}

For the packet application put $H_\delta=\Pi\cap\{\ell<\delta\}$.
At a target roof select original segments with $z^-_\sigma\ge\rho$,
$0<\ell_\sigma<\delta$, two nongrazing endpoints, and admitted target
in $H_\delta$ for the tests of \eqref{eq:v96-tests}. Define
\[
 \mathcal M_{\rho,\delta}^{\varepsilon,\rm adm}(t)
  =\int_{\text{this target family}}
    \min\{f_\sigma^\varepsilon(a_\sigma+),
           f_\sigma^\varepsilon(b_\sigma-)\}\,d\sigma.
\]
The minimum is still on the two ends of the same original segment.
The continuation anchors do not replace either endpoint.

\begin{theorem}[Decay of the recurrent two-nongrazing packet]
\label{thm:v96-packet-decay}
At each previously fixed positive entrance-gradient cutoff $\rho$,
\begin{equation}\label{eq:v96-packet-height}
 \mathcal H_M(\mathcal M_{\rho,\delta}^{\varepsilon,\rm adm})
 \le C(D/\beta)\delta^{1/32}
                        +C(1+\varepsilon^{-1})\delta.
\end{equation}
Consequently this complete admitted packet has zero ordered
zero-residence charge for fixed $\varepsilon,D,\beta$ and a fixed
comparison radius, without any lower residence cutoff. The statement
also holds for target cutoff $z(t)\ge\rho$ when
$\delta\le\rho^2/(8\Lambda_\circ)$.

If admission is tested for $H=\Pi$ instead of $H_\delta$, the
corresponding two-nongrazing packet obeys
\begin{equation}\label{eq:v96-packet-outer}
 \mathcal H_M(\mathcal M^{\varepsilon,\Pi\text{-}\rm adm}_{\rho,\delta})
 \le C(D/\beta)\{\varepsilon^{1/16}+\varepsilon^{-1}\sqrt r\}
                         +C(1+\varepsilon^{-1})\delta.
\end{equation}
This requires only returns to the same exact event, not returns to
short exact-event segments.
\end{theorem}
\begin{proof}
The two-ended variation proof of retained module 226 applies to any
Borel subfamily at the target roof. It gives
\[
 \mathcal M_{\rho,\delta}^{\varepsilon,\rm adm}(t)
 \le R_{H_\delta^{\rm adm}}^\varepsilon(t)
   +C(1+\varepsilon^{-1})
                       \int_{t-\delta}^{t+\delta}v^\rho(u)\,du.
\]
Indeed the absolute difference between a target value and its
same-segment endpoint minimum is bounded by the whole derivative
integral. The physical Borel parametrization counts that integral
once. Theorem~\ref{thm:v96-source-height} estimates the first term;
the retained positive-cutoff theorem estimates the second. This
proves \eqref{eq:v96-packet-height}. For the target-gradient version,
radial growth bounds the gradient on the entire short segment below
by $\rho/2$, as in the retained packet corollary.
For \eqref{eq:v96-packet-outer}, use the same variation comparison
with admission to $\Pi$, apply \eqref{eq:v96-weighted-anchor}, and
use the original fixed-width bounds for $b^\varepsilon$ and $p$.
\end{proof}

The new assertion concerns the positive minimum packet, including
segments with two nonzero traces. It is not the previous grazing-ended
zero theorem. It can apply when a short exact-label segment is followed
by arbitrarily many section-decision crossings, provided the same label
and same selector occupy the required total roof fraction. Neither a
bound on the number of crossings nor a pooling of their labels is used.
The packet carried by targets failing these tests is not estimated by
\eqref{eq:v96-packet-height}.
